\documentclass[conference]{IEEEtran}

\usepackage{graphicx}
\usepackage{booktabs}
\usepackage{amsmath}
\usepackage{url}
\usepackage{subcaption}
\usepackage{placeins}

\title{Multi-Task Hand Gesture Understanding :\\
From RGB Baseline to Mask-Guided }

\author{
\IEEEauthorblockN{Tianhao Liang}
\IEEEauthorblockA{
COMP0248 Coursework 1\\
}
}

\begin{document}
\maketitle

\begin{abstract}
This report presents a multi-task perception system for hand gesture understanding. The model mainly uses the idea of Multi-task, jointly learns segmentation and classification, and gets detection bounding boxes directly from predicted masks. Starting from a ResNet18 + UNet RGB baseline, I evaluate the impact of RGB-D inputs, a larger ResNet50 backbone, mask-guided classification, and staged optimization trainning strategy. After a series of comparative experiments, the final proposed architecture (RGB input + ResNet18 backbone + mask-guided + staged training) achieves the highest overall validation performance: SegIoU 0.9573, Top-1 Acc 0.9825, Macro-F1 0.9825, Detection@0.5 1.0000, and a Joint Score of 0.9699. I analyze the mechanisms that drive these improvements, and discuss current limitations and potential future work.
\end{abstract}

\section{Introduction}
Robust hand gesture perception is essential for human-robot interaction. In real-world robotic applications, systems must not only classify gestures accurately but also precisely localize and segment the hand position and outline, because these capabilities form the foundation for robots to perform various tasks, such as interacting with humans and grasping objects.

This coursework demonstrates clear model architectures, quantitative analysis, and rigorous comparisons between a baseline and improved architectures. I focus on a multi-task pipeline that outputs hand masks, gesture classes, and bounding boxes in a single forward pass, significantly enhancing inference efficiency.

Besides, I aim to address the severe ``cold-start'' and negative-transfer problems in multi-task learning, where the untrained predictions of one task may interfere with the learning of another. 

From a structural perspective, I systematically compare five configurations: (1) a standard RGB baseline; (2) a modality comparison replacing RGB with RGB-D input; (3) a backbone comparison replacing ResNet18 with ResNet50; (4) an staged train strategy without mask guidance; and (5) a final model combining mask-guided classification with staged optimization. 

\section{Dataset}
The dataset consists of 10 gesture classes : \texttt{call}, \texttt{dislike}, \texttt{like}, \texttt{ok}, \texttt{one}, \texttt{palm}, \texttt{peace}, \texttt{rock}, \texttt{stop}, and \texttt{three}.

I have 3,199 valid paired samples with a balanced class distribution. Each sample comprises a 3-channel RGB frame, a 1-channel depth map, and a binary ground-truth hand mask.

The dataset has several challenging characteristics, including high similarity in different hand gesture, captured from different viewpoints and distances, and cluttered backgrounds. These factors increase the difficulty of the task and at the same time are beneficial for training models with stronger robustness.

\section{Methodology}
\subsection{Baseline Architecture}
I adopt a shared-encoder multi-task architecture:
\begin{itemize}
    \item \textbf{Encoder}: I utilize a ResNet-18 backbone~\cite{he2016resnet}, initialized with ImageNet-1K pre-trained weights. To accommodate depth information in the RGB-D configuration, I adapt the first convolutional layer to accept 4 channels, initializing the extra depth-channel weights with the mean of the RGB-channel weightss.
    \item \textbf{Segmentation Branch}: I design a U-Net-style decoder~\cite{ronneberger2015unet} utilizing sequential upsampling blocks. Each block uses bilinear interpolation, feature concatenation with the corresponding encoder level , and two rounds of Conv-BN-ReLU layers. The final $1 \times 1$ convolution generates the output logits, which are transformed into a binary mask using a sigmoid activation and a threshold of 0.5.
    \item \textbf{Classification Branch}: The baseline applies global average pooling (GAP) to the final encoder feature map ($\mathbf{f}_5$), followed by a dropout layer ($p=0.5$) for regularization, and a linear layer mapping to the 10 gesture classes.
    \item \textbf{Detection Branch}: Rather than introducing a parameter-heavy detection head, I derive bounding boxes directly from the predicted binary segmentation mask. 
\end{itemize}

\subsection{Mask-Guided Classification}
A major challenge is that the dataset is captured within a cluttered background, while hand only ocuppies a small region. But the GAP in the baseline treats all spatial locations equally, potentially has a negative impact on classifing hand gestures because of the background information. To resolve this, I introduce a \textit{Mask-Guided Classification} head after verifing the segmentation is really precise and valid. The predicted segmentation mask ($\mathbf{m}$) is downsampled via bilinear interpolation to match the spatial resolution of the deepest encoder feature map ($\mathbf{f}_5$). I then employ spatial weighted average pooling:
\begin{equation}
\mathbf{f}_{cls} = \frac{\sum_{h,w} (\mathbf{f}_5 \odot \mathbf{mask})}{\sum_{h,w} \mathbf{mask} + \epsilon} \ ,
\end{equation}
where $\odot$ denotes element-wise multiplication and $\epsilon$ prevents division by zero. This mechanism forces the classifier to focus exclusively on the hand region, instead of being interfered by clutterd backgound. To prevent task interference, a stop-gradient operation is applied to the predicted mask prior to pooling, ensuring that the segmentation mask is trained solely under the segmentation loss, without being influenced by the classification loss.

\subsection{Loss Functions and Model Selection}
The overall training objective is a weighted sum of three tasks losses:
\begin{equation}
\mathcal{L} = \lambda_{seg}\left(\mathcal{L}_{BCE} + \mathcal{L}_{Dice}\right) + \lambda_{cls}\mathcal{L}_{CE} \ .
\end{equation}
For segmentation, I combine Binary Cross-Entropy ($\mathcal{L}_{BCE}$) for stable pixel-level gradients and Dice Loss ($\mathcal{L}_{Dice}$) to directly optimize the Intersection over Union, effectively handling foreground-background imbalance. For classification, I apply Cross-Entropy ($\mathcal{L}_{CE}$) with label smoothing (factor 0.1) to mitigate overconfidence on this relatively small dataset. Empirically, I set $\lambda_{seg} = 1.0$ and $\lambda_{cls} = 1.0$.

To eliminate subjective bias during hyperparameter tuning and checkpointing, I adopt a unified validation metric. The model state saved during training is strictly the one maximizing the Joint score:
\begin{equation}
\text{Joint} = 0.5 \times \text{SegIoU} + 0.5 \times \text{Top1Acc} \ .
\end{equation}

\subsection{Warm-Start Staged Fine-Tuning Strategy}
When attempting to train the mask-guided architecture from scratch, I observed a severe ``cold-start'' problem: in early epochs, the untrained segmentation mask acts as random noise, misguiding the classification pooling and destroying the pre-trained ImageNet features. 

To overcome this, I propose a \textit{Warm-Start Staged Fine-Tuning} strategy. First, rather than training from scratch, the model is initialized with the converged weights of the standard RGB baseline. This ensures the initial masks are already high-quality, bypassing the cold-start noise. Second, because transitioning from Global Average Pooling to mask-weighted pooling significantly alters the numerical distribution of the pooled features, the classification head must re-adapt. I freeze the entire ResNet-18 encoder (\texttt{requires\_grad=False}) for Stage 1 (Epochs 1--20). This shields the converged encoder representations from the large gradients generated during the classification head's re-adaptation phase. In Stage 2 (Epochs 21--60), the encoder is unfrozen and jointly fine-tuned at a reduced learning rate ($0.02\times$ base LR).

\section{Experimental Setup}
To ensure fair comparison, all ResNet18 experiments (baseline, RGBD, and freeze variants) adhere to an identical training protocol:
\begin{itemize}
    \item \textbf{Optimizer \& Schedule}: AdamW optimizer~\cite{loshchilov2019adamw} with a base learning rate of $5\times10^{-4}$ and weight decay of $1\times10^{-4}$. The learning rate schedule incorporates a 3-epoch linear warmup (from $0.01\times$ to $1.0\times$ base LR) followed by cosine annealing down to $1\times10^{-6}$.
    \item \textbf{Epochs \& Batch Size}: Models are trained for a maximum of 60 epochs with a batch size of 8. Early stopping is enabled, monitoring the Joint score with a patience of 20 epochs. Mixed precision (\texttt{torch.amp}) is utilized to accelerate training on NVIDIA RTX GPUs.
    \item \textbf{Data Augmentation}: An extensive augmentation pipeline via \texttt{torchvision.transforms.v2}. Geometric transformations (random rotation $\pm15^\circ$, affine scaling/translation, and resized crops) are applied synchronously to the RGB frame, depth map, and ground-truth mask to preserve spatial alignment. Modality-specific augmentations include \texttt{ColorJitter} and Gaussian blur applied exclusively to the RGB channels, alongside Gaussian noise applied independently to the depth map to simulate real-world sensor inaccuracies.
\end{itemize}

For backbone-scale comparison, I replace the ResNet18 encoder with ResNet50 in the RGBD setting (experiment \texttt{R50 RGBD tuned\_v1}). To ensure fair comparison while accounting for the larger model's increased overfitting risk, I apply conservative hyperparameter adjustments: learning rate reduced to $3\times10^{-4}$ and weight decay increased to $5\times10^{-4}$. This protocol isolates backbone capacity effects while maintaining optimization stability.

% =========================================================================
% 修复1：训练曲线图升级为双栏大图，保证清晰度，并放置在第四节末尾。
% =========================================================================
\begin{figure*}[!tp]
    \centering
    \includegraphics[width=0.9\linewidth]{figures/curves_baseline.png}
    \vspace{0.1cm}
    
    \caption{Training curves of the RGB baseline.}
    \label{fig:curve_baseline}
    
    \vspace{0.4cm}
    \includegraphics[width=0.9\linewidth]{figures/curves_best_maskguided_freeze20.png}
    \vspace{0.1cm}
    
    \caption{Training curves of the final model (mask-guided + staged training).}
    \label{fig:curve_best}
\end{figure*}


\FloatBarrier
\section{Results and Analysis}

% =========================================================================
% 修复2：主结果表，打散放在第五节开头
% =========================================================================
\begin{table*}[!t]
\centering
\caption{Performance comparison across all model configurations. All models are trained and evaluated under an identical protocol. Best results are highlighted in bold.}
\label{tab:main_results}
\begin{tabular}{lcccccc}
\toprule
Model & SegIoU & Dice & Det@0.5 & BBoxIoU & ClsAcc & MacroF1 \\
\midrule
R18 RGB baseline & 0.9492 & 0.9734 & 1.0000 & 0.9532 & 0.9750 & 0.9749 \\
R18 RGBD tuned\_v1 & 0.9478 & 0.9725 & 1.0000 & 0.9511 & 0.9775 & 0.9777 \\
R50 RGBD tuned\_v1 & 0.9386 & 0.9669 & 0.9675 & 0.9171 & 0.9400 & 0.9407 \\
R18 RGB freeze20 & 0.9544 & 0.9761 & 0.9550 & 0.9248 & 0.9450 & 0.9451 \\
\textbf{R18 RGB maskguided\_freeze20 (Ours)} & \textbf{0.9573} & \textbf{0.9780} & \textbf{1.0000} & \textbf{0.9584} & \textbf{0.9825} & \textbf{0.9825} \\
\bottomrule
\end{tabular}
\end{table*}

Table~\ref{tab:main_results} presents the quantitative performance of the five primary model configurations across segmentation, detection, and classification metrics.

\subsection{Performance of the RGB Baseline}
The standard RGB baseline establishes a remarkably robust foundation, achieving a SegIoU of 0.9492 and a Top-1 Accuracy of 0.9750. Its non-learned detection approach proves highly effective, yielding a perfect Detection Accuracy at an IoU threshold of 0.5. These results indicate that ResNet-18 features, coupled with proper data augmentation and optimization schedules, possess sufficient capacity to handle the dataset's complexity.

\subsection{Impact of Depth Information (RGB-D)}
Contrary to initial intuition, incorporating depth data (\texttt{R18 RGBD tuned\_v1}) provides only marginal gains in classification (Top-1 Acc $+0.0025$) and slightly degrades segmentation performance (SegIoU $-0.0014$). This suggests that while depth maps offer structural cues for resolving inter-class ambiguities, they may also introduce sensor noise or domain gaps. Pre-trained ImageNet filters (optimized for RGB statistics) struggle to fully exploit these 1-channel spatial maps without more sophisticated multi-modal fusion architectures.

\subsection{Backbone Scale Ablation: ResNet18 vs ResNet50}
To test whether larger model capacity improves this task, I replace the ResNet18 encoder in the RGBD setting with ResNet50 under a carefully tuned training protocol (learning rate reduced to $3\times10^{-4}$, weight decay increased to $5\times10^{-4}$ to mitigate overfitting). The result (\texttt{R50 RGBD tuned\_v1}) demonstrates a clear overfitting pattern: while achieving a training classification accuracy of 99.7\%, the validation accuracy plateaus at 94.0\%, indicating severe generalization degradation. Across all three tasks, ResNet50 underperforms ResNet18: SegIoU decreases from 0.9478 to 0.9386 ($-0.0092$), Detection@0.5 drops from 1.0000 to 0.9675 ($-0.0325$), and Top-1 accuracy drops from 0.9775 to 0.9400 ($-0.0375$). This indicates that, for the current dataset size (3,199 samples), increasing backbone capacity from 11M to 25M parameters primarily increases optimization difficulty and overfitting risk rather than improving representation quality. The larger model's capacity exceeds what the limited training data can effectively regularize, leading to task-specific feature degradation despite architectural advantages.

% =========================================================================
% 修复3：定性效果图，缩小尺寸，防止太大，插在 C 小节附近
% =========================================================================
\begin{figure*}[!t]
\centering
\begin{subfigure}{0.45\textwidth}
    \centering
    \includegraphics[width=0.85\linewidth]{figures/qual_baseline.png}
    \caption{RGB baseline}
\end{subfigure}
\hfill
\begin{subfigure}{0.45\textwidth}
    \centering
    \includegraphics[width=0.85\linewidth]{figures/qual_rgbd_tuned_v1.png}
    \caption{RGBD tuned\_v1}
\end{subfigure}

\vspace{0.3cm}

\begin{subfigure}{0.45\textwidth}
    \centering
    \includegraphics[width=0.85\linewidth]{figures/qual_freeze20.png}
    \caption{RGB freeze20}
\end{subfigure}
\hfill
\begin{subfigure}{0.45\textwidth}
    \centering
    \includegraphics[width=0.85\linewidth]{figures/qual_best_maskguided_freeze20.png}
    \caption{Best model (Ours)}
\end{subfigure}
\caption{Qualitative segmentation and bounding box examples for four primary reported configurations. The ResNet50 ablation is reported quantitatively in Table~\ref{tab:main_results} due page-limit constraints.}
\label{fig:qualitative}
\end{figure*}

\subsection{Ablation: Staged Optimization Without Mask-Guidance}
To isolate the effect of staged optimization, I apply the freeze-then-unfreeze strategy to the baseline architecture without mask-guided classification (\texttt{R18 RGB freeze20}). This experiment uses the same staged training protocol (encoder frozen for epochs 1--20, then unfrozen with reduced learning rate) but retains the standard GAP-based classification head. The result demonstrates a significant improvement in segmentation (SegIoU rises to 0.9544), confirming that staged optimization benefits dense prediction tasks. However, this comes at a severe cost to classification (Top-1 Acc drops to 0.9450) and detection stability. By freezing the backbone early on without mask-guidance, the network becomes heavily biased towards dense spatial prediction (segmentation), preventing the encoder from refining the fine-grained, abstract discriminative features necessary for accurate classification. This ablation establishes that staged optimization alone is insufficient; it must be combined with mask-guided classification to achieve balanced multi-task performance.


\subsection{The Final Model: Synergy of Mask-Guidance and Staging}
The final proposed model (\texttt{R18 RGB maskguided\_freeze20}) successfully resolves the trade-offs observed in the ablations. As demonstrated in Table~\ref{tab:joint_delta}, it achieves the highest overall Joint score, outperforming the baseline by a clear margin of $+0.0078$. 

% =========================================================================
% 修复4：Delta 表格，放单栏
% =========================================================================
\begin{table}[tb]
\centering
\caption{Joint score and absolute improvement ($\Delta$) relative to the RGB baseline.}
\label{tab:joint_delta}
\begin{tabular}{lcc}
\toprule
Model & Joint Score & $\Delta$ vs Baseline \\
\midrule
R18 RGB baseline & 0.9621 & 0.0000 \\
R18 RGBD tuned\_v1 & 0.9626 & +0.0005 \\
R50 RGBD tuned\_v1 & 0.9393 & -0.0228 \\
R18 RGB freeze20 & 0.9497 & -0.0124 \\
\textbf{R18 RGB maskguided\_freeze20} & \textbf{0.9699} & \textbf{+0.0078} \\
\bottomrule
\end{tabular}
\end{table}

By combining mask-guided spatial pooling with the freeze-then-unfreeze training paradigm, the model forces the classifier to focus exclusively on the foreground hand while preventing early inaccurate masks from corrupting the encoder. This synergy establishes state-of-the-art performance across all metrics: SegIoU (0.9573), Top-1 Accuracy (0.9825), and Mean BBox IoU (0.9584).

Figure~\ref{fig:qualitative} visualizes qualitative predictions of four primary models (excluding the ResNet50 ablation for space). While all models locate the core hand area well, the final mask-guided model produces noticeably sharper boundaries and fewer false-positive artifacts in the background, directly contributing to its superior BBox IoU. Furthermore, the confusion matrices (Figure \ref{fig:cm_comparison}) demonstrate that the final model effectively mitigates difficult inter-class confusions, such as distinguishing between the geometrically similar \texttt{ok} and \texttt{one} gestures.

% =========================================================================
% 修复5：混淆矩阵，调小宽度比例，插在第六节前面
% =========================================================================
\begin{figure*}[!t]
    \centering
    \begin{subfigure}{0.45\linewidth}
        \centering
        \includegraphics[width=0.75\linewidth]{figures/cm_baseline.png}
        \caption{RGB baseline}
        \label{fig:cm_baseline}
    \end{subfigure}
    \hfill
    \begin{subfigure}{0.45\linewidth}
        \centering
        \includegraphics[width=0.75\linewidth]{figures/cm_best_maskguided_freeze20.png}
        \caption{Final model}
        \label{fig:cm_best}
    \end{subfigure}
    \caption{Confusion matrix comparison between the baseline and our final model, illustrating resolved inter-class ambiguities.}
    \label{fig:cm_comparison}
\end{figure*}


\FloatBarrier
\section{Innovation and Impact}
Compared to a standard multi-task baseline, our proposed pipeline introduces two key innovations that cumulatively enhance perception performance:
\begin{itemize}
    \item \textbf{Mask-Guided Classification Head}: By replacing standard GAP with a mask-weighted spatial pooling mechanism, the model is compelled to explicitly associate gesture predictions with the segmented hand region, effectively suppressing background noise.
    \item \textbf{Warm-Start Staged Fine-Tuning}: The cold-start degradation inherent to training mask-guided architectures from scratch is resolved by initializing from a baseline checkpoint and freezing the encoder. This shields the backbone from classification re-adaptation gradients and ensures training stability.
\end{itemize}
The cumulative gain is significant and consistent across all evaluated metrics: segmentation precision, classification accuracy, and bounding box quality all demonstrate measurable superiority over the baseline.

\FloatBarrier
\section{Discussion and Limitations}
This series of controlled experiments reveals three critical insights into multi-task perception for robotics. First, naive optimization adjustments—such as merely freezing the encoder—can over-bias the network toward dense spatial features (aiding segmentation) while severely impairing its ability to extract abstract semantic features (harming classification). Second, while mask-guided pooling is theoretically elegant, it suffers from a severe cold-start problem when trained from scratch due to negative transfer in early epochs. Third, achieving state-of-the-art performance necessitates aligning architectural innovations with tailored optimization strategies; the complementary effects of mask-guidance and warm-start staged fine-tuning unlock the model's full potential.

Despite these advancements, several limitations persist. First, given the relatively small subject pool (33 individuals), the reported performance margins between model configurations could be susceptible to single-run variance, and multi-seed evaluations would further strengthen the statistical reliability of these conclusions. Second, relying on a thresholded binary mask to gate the classification head renders the model susceptible in scenarios involving extreme motion blur or occlusion, where the mask prediction might momentarily fail. 

Future work should explore multi-seed confidence intervals and incorporate temporal modeling (e.g., via 3D convolutions or LSTMs) to enforce video-level consistency across adjacent frames. Furthermore, investigating lightweight deployment optimizations (e.g., TensorRT or ONNX conversion) remains vital for real-time inference on edge robotic platforms.

\section{Conclusion}
This report establishes a robust, reproducible multi-task pipeline for simultaneous hand gesture segmentation, classification, and detection. I systematically evaluated the baseline architecture against RGB-D input modalities, a larger ResNet50 backbone, optimization ablations, and structural modifications. The ResNet50 control experiment reveals a critical insight: despite careful hyperparameter tuning to mitigate overfitting, the larger model (25M vs 11M parameters) exhibits severe generalization degradation, with training classification accuracy reaching 99.7\% while validation accuracy plateaus at 94.0\%. This confirms that larger capacity does not guarantee better performance under the current data scale (3,199 samples), and that model selection must balance representational power with dataset-appropriate regularization. Our proposed final model, \texttt{R18 RGB maskguided\_freeze20}, successfully integrates mask-guided spatial pooling with a staged freeze-then-unfreeze optimization strategy. This approach decisively overcomes the cold-start problem in multi-task learning, suppresses background interference, and achieves superior comprehensive performance across all key metrics. Consequently, it demonstrates exceptional practical viability for real-world human-robot interaction systems.

\bibliographystyle{IEEEtran}
\bibliography{references}

\end{document}